% Generated draft: requires mathematical and source review.
\documentclass[11pt]{article}
\usepackage[margin=25mm]{geometry}
\usepackage{fontspec}
\setmainfont{Latin Modern Roman}
\usepackage{amsmath,amssymb,mathtools,unicode-math}
\setmathfont{Latin Modern Math}
\usepackage{enumitem,needspace,xurl,hyperref}
\setlength{\parindent}{0pt}
\setlength{\parskip}{5pt}
\setlength{\emergencystretch}{3em}
\newcommand{\N}{\mathbb{N}}
\newcommand{\Z}{\mathbb{Z}}
\newcommand{\Q}{\mathbb{Q}}
\newcommand{\R}{\mathbb{R}}
\newcommand{\C}{\mathbb{C}}
\newcommand{\F}{\mathbb{F}}
\newcommand{\A}{\mathbb{A}}
\newcommand{\PP}{\mathbb{P}}
\newcommand{\E}{\mathbb{E}}
\newcommand{\eps}{\varepsilon}
\newcommand{\dd}{\,\mathrm{d}}
\begin{document}
\section*{UnsolvedMath ID: 20003282}
\textbf{Source problem code:} catalogue entry
\subsection{Definitions and mathematical statement}
The source asks, with attribution to G. Margulis [130], for equidistribution of $u(t_n)x$ in $G/\Gamma$ for three sampling families: a damaged power/floor expression printed as $[n\alpha]3$ with $\alpha>1$; $t_n=\lfloor P(n)\rfloor$ for a polynomial $P$; and the $n$th prime. The aid says the first expression has a misleading rendering artifact, but the supplied record does not reproduce its full repaired formula or the preceding Section 4.1 hypotheses. The placement of the power of $n$ and the trailing $3$ therefore remain unresolved here. Its being an integer-valued floor sequence is sufficient for the first literal obstruction below. The source also omits the target measure, orbit-closure hypotheses and whether $P$ is nonconstant, integer-valued or real-coefficient.

For a specified probability measure $m$ on the relevant space or orbit closure, equidistribution means weak convergence
\[
 \frac1N\sum_{n=1}^N\delta_{u(t_n)x}\ \longrightarrow\ m,
\]
with the appropriate continuous-test-function convention. On a noncompact quotient escape of mass and finite-volume hypotheses need attention; the source provides none locally.

The aid's literal counterexample takes $G=\R$, $\Gamma=\Z$, $u(t)=t$ acting by addition on $\R/\Z$. Every continuous-time orbit is the entire circle, but every integer sampling time gives $u(t_n)x=x$, so the empirical measure is $\delta_x$, not Haar. Its applicability to the original workshop setting depends on the missing Lie-group/unipotence hypotheses actually allowing this abelian example; no stronger semisimple formulation is declared refuted merely from this local extraction.

A sharper polynomial obstruction in the aid keeps time one uniquely ergodic. Choose irrational $0<\omega<1$, let $u(t)y=y+\omega t$ on $\R/\Z$, and use the real polynomial $P(s)=s/\omega$. Then
\[
 u(\lfloor P(n)\rfloor)x
 =x-\omega\{n/\omega\}\pmod1.
\]
The reported limiting measure is the pushforward of Lebesgue measure on $[0,1)$ by $r\mapsto x-\omega r$, hence density $1/\omega$ on an arc of length $\omega$, rather than Haar on the circle. Thus unique ergodicity of the time-one map alone does not repair the unrestricted real-polynomial statement.

The complementary positive result reported by the aid concerns totally irrational torus translations: for $u(t)y=y+t v$ on $\R^d/\Z^d$, require $k\cdot v\notin\Q$ for every $0\ne k\in\Z^d$, equivalently rational independence of $1,v_1,\ldots,v_d$. It reports the intended power-family and prime equidistribution, and the polynomial case when $P$ is nonconstant and integer-valued. For real-coefficient $P$ it gives a sufficient Fourier nonresonance condition called (NR), but no formula for that condition is supplied. These torus conclusions are not extended to arbitrary homogeneous dynamics, and the first positive family needs the repaired time formula before a precise theorem can be asserted.

The source attributes a proposed proof for case (1) with $\alpha$ close to one to A. Venkatesh and ends with an incomplete sentence about Ces\`aro averages. These are dated/truncated context, not an established proof or a recovered theorem. The registry's refuted/solved-in-literature labels are retained as metadata for the literal readings, not a blanket resolution of all intended sparse-unipotent questions.
\subsection{Short English statement}
Integer sampling can miss continuous-time equidistribution, and a real floored polynomial can defeat even an irrational time-one rotation. The aid reports positive torus results under stronger arithmetic hypotheses, but the first time formula, general group setup and nonresonance condition are incomplete in the supplied record.
\subsection{Sources}
\begin{itemize}
\item Local UnsolvedMath registry, permanent ID 20003282, source code catalogue entry; \texttt{aim-other-notes.json}, notes[88], source tag question.
\item AIM, \emph{Emerging applications of measure rigidity}, Question 15, G. Margulis [130]: \url{https://aimath.org/WWN/measrigid/measrigid.pdf}. The supplied verification says the PDF and Section 4.1 were checked but does not reproduce the complete setup; no [130] entry is supplied.
\item A. Venkatesh's suggested near-$1$ exponent proof is mentioned in the original text. The Ces\`aro sentence is truncated. The positive torus and (NR) results have no full external references in the summary.
\item Machine-generated research attempt 1: counterexample, refuted at the run, rigor labelled ``proved,'' low novelty confidence; registry solved/``SOLVED-IN-LITERATURE'' and reconstructed unverified statement. Not peer reviewed; independently verify before citation. Registry updated 14 August 2026.
\end{itemize}
\end{document}
